% Bibliography — Hyperbit, Landauer, FFGFT Corpus

\noindent\textbf{Matzke Sources}

\begin{description}[leftmargin=3.5em,style=sameline,itemsep=0.4em]

\item[{[M1]}] D.~Matzke, \textit{Black Holes from Hyperbits},
  IPI Annual Conference, August 2026. \\
  Main source for the hyperbit framework; derivation of Schwarzschild radius
  and Bekenstein entropy from $e_i^2 = +1$.

\item[{[M2]}] D.~Matzke, \textit{G(6,\,$\mathbb{Z}_3\text{C}$) Hyperspin
  Framework}, IPI Working Paper, 2026. \\
  Algebraic foundation; $\mathrm{Cl}(6)$ and fermion generations.
  Compared with FFGFT in Doc.\,324
  \href{https://github.com/jpascher/T0-Time-Mass-Duality/blob/main/2/pdf/324_G6_Z3C_FFGFT_Vergleich_De.pdf}{(PDF)}.

\end{description}

\bigskip
\noindent\textbf{FFGFT Corpus (J.\,Pascher)}

\begin{description}[leftmargin=3.5em,style=sameline,itemsep=0.4em]

\item[{[P1]}] J.~Pascher, \textit{A271 — Landauer and Phase-Space Bookkeeping},
  FFGFT Corpus, July 2026.\\
  \href{https://github.com/jpascher/T0-Time-Mass-Duality/blob/main/A_Serie_Export/pdf/A271_Landauer_De.pdf}{Direct Link PDF}
  $\cdot$ Core source for Ch.\,\ref{ch:landauer} and \ref{ch:landauer_phasenraum}.

\item[{[P2]}] J.~Pascher, \textit{A272 — Carrier and Information},
  FFGFT Corpus, July 2026.\\
  \href{https://github.com/jpascher/T0-Time-Mass-Duality/blob/main/A_Serie_Export/pdf/A272_Traeger_und_Information_De.pdf}{Direct Link PDF}
  $\cdot$ Textual criticism of the Landauer reception; Ch.\,\ref{ch:traeger}.

\item[{[P3]}] J.~Pascher, \textit{A273 — The Computing Sphere},
  FFGFT Corpus, July 2026.\\
  \href{https://github.com/jpascher/T0-Time-Mass-Duality/blob/main/A_Serie_Export/pdf/A273_Die_Rechenkugel_De.pdf}{Direct Link PDF}
  $\cdot$ Abacus, colloid, limiting transition; Ch.\,\ref{ch:rechenkugel}.

\item[{[P4]}] J.~Pascher, \textit{A280/A281 — Landauer Phase-Space Bookkeeping
  (Extended) / Carrier and Information (Extended)},
  FFGFT Corpus, 2026.\\
  \href{https://github.com/jpascher/T0-Time-Mass-Duality/blob/main/A_Serie_Export/pdf/A280_Landauer_Phasenraum_Buchhaltung_De.pdf}{A280 PDF}
  $\cdot$
  \href{https://github.com/jpascher/T0-Time-Mass-Duality/blob/main/A_Serie_Export/pdf/A281_Traeger_und_Information_De.pdf}{A281 PDF}.

\item[{[P5]}] J.~Pascher, \textit{A283 — Thermodynamics and FFGFT},
  FFGFT Corpus, 2026.\\
  \href{https://github.com/jpascher/T0-Time-Mass-Duality/blob/main/A_Serie_Export/pdf/A283_Thermodynamik_FFGFT_De.pdf}{Direct Link PDF}
  $\cdot$ Ch.\,\ref{ch:thermodynamik}.

\item[{[P6]}] J.~Pascher, \textit{A284 — Measurement and Projection},
  FFGFT Corpus, 2026.\\
  \href{https://github.com/jpascher/T0-Time-Mass-Duality/blob/main/A_Serie_Export/pdf/A284_Messung_Projektion_De.pdf}{Direct Link PDF}
  $\cdot$ Ch.\,\ref{ch:messung}.

\item[{[P7]}] J.~Pascher, \textit{Doc.\,257 — Information Unit on $T^4$},
  FFGFT Corpus.\\
  \href{https://github.com/jpascher/T0-Time-Mass-Duality/blob/main/2/pdf/257_Informationseinheit_Skala_De.pdf}{Direct Link PDF}
  $\cdot$ $\Ebit = \hbar c/L$, winding quantum $\Delta w = 1$; Ch.\,\ref{ch:informationseinheit_skala}.

\item[{[P8]}] J.~Pascher, \textit{Doc.\,184 — The p-Bit},
  FFGFT Corpus.\\
  \href{https://github.com/jpascher/T0-Time-Mass-Duality/blob/main/2/pdf/184_p-bit_De.pdf}{Direct Link PDF}
  $\cdot$ $L_8 \approx 22\,\text{nm}$, Stoner-Wohlfarth; Ch.\,5.

\item[{[P9]}] J.~Pascher, \textit{Doc.\,326 — Matzke--FFGFT Comparison},
  FFGFT Corpus, August 2026.\\
  \href{https://github.com/jpascher/T0-Time-Mass-Duality/blob/main/2/pdf/326_Matzke_FFGFT_Vergleich_De.pdf}{Direct Link PDF}
  $\cdot$ Basis for Ch.\,10.

\item[{[P10]}] J.~Pascher, \textit{Doc.\,324 — G(6,\,$\mathbb{Z}_3$C)--FFGFT Comparison},
  FFGFT Corpus, 2026.\\
  \href{https://github.com/jpascher/T0-Time-Mass-Duality/blob/main/2/pdf/324_G6_Z3C_FFGFT_Vergleich_De.pdf}{Direct Link PDF}.

\item[{[P11]}] J.~Pascher, \textit{Doc.\,325 — Hawking and FFGFT},
  FFGFT Corpus, 2026.\\
  \href{https://github.com/jpascher/T0-Time-Mass-Duality/blob/main/2/pdf/325_Hawking_FFGFT_De.pdf}{Direct Link PDF}.

\item[{[P12]}] J.~Pascher, \textit{Doc.\,302 — Unit Cells and Bit Capacity},
  FFGFT Corpus, 2026.\\
  \href{https://github.com/jpascher/T0-Time-Mass-Duality/blob/main/2/pdf/302_Elementarzellen_Spektrale_Antwort_De.pdf}{Direct Link PDF}.

\item[{[P13]}] J.~Pascher, \textit{Doc.\,301 — Information Reversal Test},
  FFGFT Corpus, 2026.\\
  \href{https://github.com/jpascher/T0-Time-Mass-Duality/blob/main/2/pdf/301_Information_Umkehrtest_De.pdf}{Direct Link PDF}.

\item[{[P14]}] J.~Pascher, \textit{Doc.\,174 — Spin and Qubit States},
  FFGFT Corpus.\\
  \href{https://github.com/jpascher/T0-Time-Mass-Duality/blob/main/2/pdf/174_T0_Spin_Qubit_Zustaende_De.pdf}{Direct Link PDF}.

\end{description}

\bigskip
\noindent\textbf{Mathematical Foundations (External Sources)}

\begin{description}[leftmargin=3.5em,style=sameline,itemsep=0.4em]

\item[{[T1]}] A.~Hatcher, \textit{Algebraic Topology},
  Cambridge University Press, 2002.\\
  Freely available: \href{https://pi.math.cornell.edu/~hatcher/AT/AT.pdf}{pi.math.cornell.edu/\textasciitilde hatcher/AT/AT.pdf}\\
  Ch.\,1.2: Fundamental group of the circle $\pi_1(S^1) \cong \mathbb{Z}$;
  Ch.\,1.2 Example 1.22: Product formula $\pi_1(X \times Y) \cong \pi_1(X) \times \pi_1(Y)$,
  from which $\pi_1(T^n) \cong \mathbb{Z}^n$ follows.

\item[{[T2]}] G.~E.~Bredon, \textit{Topology and Geometry},
  Springer GTM 139, 1993.\\
  Ch.\,III.§5: Homotopy groups of product spaces.

\item[{[T3]}] J.~R.~Munkres, \textit{Topology}, 2nd ed.,
  Prentice Hall, 2000.\\
  Ch.\,9: Fundamental group; Ch.\,11: Seifert-van Kampen theorem.

\end{description}

\begin{description}[leftmargin=3.5em,style=sameline,itemsep=0.4em]

\item[{[V1]}] M.\,M.~Vopson,
  \textit{The mass-energy-information equivalence principle},
  AIP Advances \textbf{9}(9), 095206 (2019).\\
  Original formulation of the mass-information equivalence;
  basis for all Vopson discussions in Ch.\,3, 4, 9.

\item[{[V2]}] M.\,M.~Vopson,
  \textit{The information catastrophe},
  AIP Advances \textbf{10}(8), 085014 (2020).\\
  Estimation of the bit capacity of the universe;
  $\sim 10^{80}$ to $10^{90}$ bits.

\item[{[V3]}] M.\,M.~Vopson,
  \textit{Second law of infodynamics and its implications
  for the simulated universe hypothesis},
  AIP Advances \textbf{13}, 105308 (2023).\\
  2nd law of infodynamics; symmetry minimization as an
  information principle (Hund's rules).

\item[{[V4]}] M.\,M.~Vopson,
  \textit{Experimental protocol to test the second law of
  infodynamics using a scanning tunnelling microscope},
  AIP Advances \textbf{14}, 035205 (2024).\\
  Experimental proposal for detecting infodynamics;
  dark matter as stored information.

\end{description}

\begin{description}[leftmargin=3.5em,style=sameline,itemsep=0.4em]

\item[{[B1]}] R.~Landauer, \textit{Irreversibility and Heat Generation in the
  Computing Process}, IBM J.\,Res.\,Dev.\ \textbf{5}, 183 (1961). \\
  Primary source; ``logically irreversible'' is the precise formulation.

\item[{[B2]}] Liouville theorem: e.g.\ Landau/Lifshitz,
  \textit{Mechanics}, §46.

\item[{[B3]}] C.~H.~Bennett, \textit{The Thermodynamics of Computation},
  Int.\ J.\ Theor.\ Phys.\ \textbf{21}, 905 (1982); idem,
  \textit{Notes on Landauer's Principle},
  Stud.\ Hist.\ Phil.\ Mod.\ Phys.\ \textbf{34}, 501 (2003).

\item[{[B4]}] A.~Bérut et al., \textit{Experimental verification of
  Landauer's principle}, Nature \textbf{483}, 187 (2012). \\
  Colloid experiment; directly confirms the Landauer limit.

\item[{[B5]}] Y.~Jun, M.~Gavrilov, J.~Bechhoefer,
  \textit{High-Precision Test of Landauer's Principle in a Feedback Trap},
  Phys.\ Rev.\ Lett.\ \textbf{113}, 190601 (2014).

\item[{[B6]}] C.~E.~Shannon, \textit{A Mathematical Theory of Communication},
  Bell System Technical Journal \textbf{27}, 379 (1948). \\
  Information-theoretic bit ($\log_2$), not thermodynamic.

\item[{[B7]}] E.~T.~Jaynes, \textit{Information Theory and Statistical Mechanics},
  Phys.\ Rev.\ \textbf{106}, 620 (1957). \\
  Shannon-Boltzmann connection; applies exactly only for uniform distribution.

\item[{[B8]}] C.~Furey, \textit{Standard Model Physics from an Algebra?},
  Diss., Waterloo, 2015. \\
  $\mathrm{Cl}(6)$ and fermion generations; foundation of \textbf{[M1/M2]}.

\item[{[B9]}] J.~Bekenstein, \textit{Black holes and entropy},
  Phys.\ Rev.\ D \textbf{7}, 2333 (1973). \\
  Entropy-area relation; one of the three bridges in \textbf{[M1]}.

\item[{[B10]}] T.~Sagawa, M.~Ueda,
  \textit{Second Law of Thermodynamics with Discrete Quantum Feedback Control},
  Phys.\ Rev.\ Lett.\ \textbf{100}, 080403 (2008).

\item[{[B11]}] E.~Fredkin, T.~Toffoli,
  \textit{Conservative Logic},
  Int.\ J.\ Theor.\ Phys.\ \textbf{21}, 219 (1982). \\
  Reversible gates; Landauer-cost-free computing is in principle possible.

\item[{[B12]}] W.~K.~Wootters, W.~H.~Zurek,
  \textit{A single quantum cannot be cloned},
  Nature \textbf{299}, 802 (1982). No-cloning theorem.

\item[{[B13]}] P.~W.~Shor, Phys.\ Rev.\ A \textbf{52}, R2493 (1995);
  A.~M.~Steane, Phys.\ Rev.\ Lett.\ \textbf{77}, 793 (1996). \\
  Quantum error correction.

\end{description}

\bigskip
\noindent{\small\textcopyright\ 2026 Johann Pascher
$\cdot$ \href{https://creativecommons.org/licenses/by/4.0/}{CC BY 4.0}
$\cdot$ ORCID \href{https://orcid.org/0009-0000-6518-4064}{0009-0000-6518-4064}}